\setchapterpreamble[u]{\margintoc}
\chapter{The Singular Value Decomposition}
\labch{svd}

\sources{Strang, \emph{Introduction to Linear Algebra}, 6th ed.\ \cite{strang2023ila},
§7.1--7.4; MIT 18.06 \cite{ocw1806}, Lecture~31.}

Eigenvalues require a square matrix and may fail to supply a full basis. The
singular value decomposition has neither restriction: every matrix has one. It
is the factorisation that ties the four subspaces together and the one that
underlies most of \refch{learning-from-data}.

\section{The statement}
\labsec{svd-statement}

\begin{theorem}
Every \(A\in\R^{m\times n}\) can be written
\[
	A=U\Sigma V\tp,
\]
where \(U\in\R^{m\times m}\) and \(V\in\R^{n\times n}\) are orthogonal and
\(\Sigma\in\R^{m\times n}\) is diagonal with entries
\(\sigma_1\ge\sigma_2\ge\cdots\ge\sigma_r>0\) and zeros elsewhere, \(r\) being
the rank of \(A\).
\end{theorem}

The columns \(\vect{v}_j\) are the \emph{right singular vectors}, the columns
\(\vect{u}_i\) the \emph{left singular vectors}, and the \(\sigma_i\) the
\emph{singular values}. The content of the factorisation is the pairing
\[
	A\vect{v}_i=\sigma_i\vect{u}_i,
	\qquad i=1,\ldots,r,
\]
which says that \(A\) maps one orthonormal set to another, up to individual
stretching factors.

\marginnote{Compare with eigenvectors, where the same basis appears on both
sides. Allowing two different bases is exactly what removes the need for \(A\)
to be square or diagonalisable.}

\section{Where it comes from}
\labsec{svd-construction}

Suppose \(A=U\Sigma V\tp\). Then
\[
	A\tp A=V\Sigma\tp U\tp U\Sigma V\tp=V(\Sigma\tp\Sigma)V\tp,
	\qquad
	AA\tp=U(\Sigma\Sigma\tp)U\tp .
\]
Both are symmetric and positive semidefinite, so \refsec{spectral} applies. The
\(\vect{v}_i\) are eigenvectors of \(A\tp A\), the \(\vect{u}_i\) are
eigenvectors of \(AA\tp\), and the singular values are the square roots of the
shared non-zero eigenvalues.

\begin{example}
Let \(A=\begin{bmatrix}3&0\\4&5\end{bmatrix}\). Then
\(A\tp A=\begin{bmatrix}25&20\\20&25\end{bmatrix}\), with eigenvalues \(45\) and
\(5\) and eigenvectors \([1,1]\tp\) and \([1,-1]\tp\). Hence
\(\sigma_1=3\sqrt5\), \(\sigma_2=\sqrt5\), and
\[
	\vect{v}_1=\tfrac{1}{\sqrt2}\begin{bmatrix}1\\1\end{bmatrix},
	\quad
	\vect{u}_1=\frac{A\vect{v}_1}{\sigma_1}=\tfrac{1}{\sqrt{10}}\begin{bmatrix}1\\3\end{bmatrix},
	\quad
	\vect{u}_2=\frac{A\vect{v}_2}{\sigma_2}=\tfrac{1}{\sqrt{10}}\begin{bmatrix}3\\-1\end{bmatrix}.
\]
As a check, \(\vect{u}_1\tp\vect{u}_2=0\) and
\(\sigma_1\sigma_2=3\sqrt5\cdot\sqrt5=15=\det A\).
\end{example}

\begin{remark}
In floating point one does not form \(A\tp A\); squaring the matrix squares its
condition number. The derivation above explains what the singular vectors are;
production algorithms compute them by orthogonal transformations applied to
\(A\) directly.
\end{remark}

\section{The four subspaces, one picture}
\labsec{svd-subspaces}

Split the singular vectors at index \(r\):
\[
	\underbrace{\vect{v}_1,\ldots,\vect{v}_r}_{\text{row space}},\quad
	\underbrace{\vect{v}_{r+1},\ldots,\vect{v}_n}_{\Null(A)},\qquad
	\underbrace{\vect{u}_1,\ldots,\vect{u}_r}_{\Col(A)},\quad
	\underbrace{\vect{u}_{r+1},\ldots,\vect{u}_m}_{\Null(A\tp)} .
\]
Each group is an orthonormal basis of the subspace beneath it. The orthogonality
statements of \refch{orthogonality} are then read off directly, and the action
of \(A\) becomes: discard the nullspace component, and stretch the surviving
\(r\) directions by \(\sigma_1,\ldots,\sigma_r\).

\section{Rank-one pieces and best approximation}
\labsec{eckart-young}

The column-times-row reading of \refsec{four-ways} turns the factorisation into
a sum ordered by importance:
\begin{equation}
	A=\sum_{i=1}^{r}\sigma_i\,\vect{u}_i\vect{v}_i\tp .
	\labeq{svd-sum}
\end{equation}

Truncating after \(k\) terms gives \(A_k\), and no other matrix of rank \(k\)
comes closer.

\begin{theorem}[Eckart--Young]
Among all matrices \(B\) of rank at most \(k\), the truncation \(A_k\) minimises
\(\lVert A-B\rVert\) in both the spectral and the Frobenius norm, with
\(\lVert A-A_k\rVert_2=\sigma_{k+1}\).
\end{theorem}

\marginnote{This single theorem is the justification for principal component
analysis, low-rank compression of weight matrices, and latent-factor models
alike. They differ in what the matrix contains, not in the mathematics applied
to it.}

\section{The pseudoinverse}
\labsec{pseudoinverse}

When \(A\) is not invertible, the closest available substitute is
\[
	A^{+}=V\Sigma^{+}U\tp,
\]
where \(\Sigma^{+}\) transposes \(\Sigma\) and replaces each \(\sigma_i>0\) by
\(1/\sigma_i\), leaving the zeros alone. Then \(\vect{x}=A^{+}\vect{b}\) is the
least-squares solution of smallest length: among all vectors minimising
\(\lVert A\vect{x}-\vect{b}\rVert\), it is the one lying entirely in the row
space.

\begin{remark}
This closes a loop opened in \refch{orthogonality}. There the normal equations
needed independent columns so that \(A\tp A\) could be inverted. The
pseudoinverse removes that condition and picks a canonical answer from the
infinitely many that rank deficiency allows.
\end{remark}
